% BHU PYQ -- Manually transcribed & error-corrected
% Paper: BCH-216 : Business Statistics
% Exam: B.Com. (Hons.) Semester III Examination, 2016-17

\documentclass[11pt,a4paper]{article}
\usepackage[a4paper,top=2.5cm,bottom=2.5cm,left=2.8cm,right=2.8cm,headheight=14pt]{geometry}
\usepackage{amsmath,amssymb}
\usepackage[T1]{fontenc}
\usepackage[utf8]{inputenc}
\usepackage{lmodern,microtype,fancyhdr,enumitem,hyperref,lastpage,parskip,booktabs}

\hypersetup{
    colorlinks=true,
    urlcolor=black,
    linkcolor=black,
    pdftitle={BCH-216: Business Statistics -- BHU B.Com. (Hons.) Sem III 2016-17}
}

\pagestyle{fancy}
\fancyhf{}
\renewcommand{\headrulewidth}{0.4pt}
\renewcommand{\footrulewidth}{0.4pt}
\fancyhead[L]{\small   \textit{Banaras Hindu University}}
\fancyhead[R]{\small   \textit{BCH-216: Business Statistics}}
\fancyfoot[C]{\small Page  \thepage\ of \pageref{LastPage}}


\newlist{parts}{enumerate}{2}
\setlist[parts,1]{label=(\alph*),leftmargin=*,itemsep=5pt,topsep=3pt}
\setlist[parts,2]{label=(\roman*),leftmargin=*,itemsep=3pt,topsep=2pt}

\newcommand{\pts}[1]{\hfill   \textit{[#1]}}


\begin{document}


\begin{center}
  {\large\bfseries BANARAS HINDU UNIVERSITY}\\[2pt]
  {
\normalsize B.Com. (Hons.) --- Semester III Examination, 2016-17}\\[6pt]
  \rule{0.8\linewidth}{0.5pt}\\[6pt]
  {\large\bfseries Paper No. BCH-216: Business Statistics}\\[4pt]
  {
\normalsize Time: 3 Hours\hfill Maximum Marks: 70}
\end{center}

\rule{\linewidth}{0.4pt}

\smallskip



\noindent   \textit{Instructions: Answer all questions. All questions carry equal marks (14 marks each).}

\bigskip




\noindent   \textbf{Question 1.}
Discuss the uses of statistics for economic analysis and planning with special reference to India.\pts{14}

\centerline{   \textbf{OR}}

Discuss the Questionnaire and schedule methods of data collection along with their merits and limitations.\pts{14}

\medskip\rule{\linewidth}{0.2pt}




\noindent   \textbf{Question 2.}
From the following data, calculate the range of marks obtained by middle 50\% of the students:\pts{14}

\begin{center}
  
\begin{tabular}{lcccccc}
     oprule
    Marks (More than) & 0 & 10 & 20 & 30 & 40 & 50 \\
    \midrule
    No. of Students & 100 & 92 & 80 & 40 & 20 & 6 \\
    ottomrule
  \end{tabular}
\end{center}

\centerline{   \textbf{OR}}

A train runs 25 kms.\ at a speed of 30 kms.\ per hour, another 50 kms.\ at a speed of 40 kms.\ per hour, then due to repair of track travels for 6 minutes at a speed of 10 kms.\ per hour and finally covers 24 kms.\ at a speed of 24 kms.\ per hour. What is the average speed (in kms.) per hour?\pts{14}

\medskip\rule{\linewidth}{0.2pt}




\noindent   \textbf{Question 3.}
A purchasing agent obtained samples of incandescent lamps from two suppliers. He had the samples tested in his own laboratory for the length of life with the following results:

\begin{center}
  
\begin{tabular}{lcc}
     oprule
    Length of Life in Hours & \multicolumn{2}{c}{Samples From} \\
    \cmidrule{2-3}
    & Company A & Company B \\
    \midrule
    700 and under 900 & 10 & 3 \\
    900 and under 1100 & 16 & 42 \\
    1100 and under 1300 & 26 & 12 \\
    1300 and under 1500 & 8 & 3 \\
    ottomrule
  \end{tabular}
\end{center}
Which company's lamps are more uniform?\pts{14}

\centerline{   \textbf{OR}}

What do you understand by skewness? Point out their role in analyzing a frequency distribution.\pts{14}

\medskip\rule{\linewidth}{0.2pt}




\noindent   \textbf{Question 4.}
What is Spearman's Rank Correlation Co-efficient? Bring out its usefulness.\pts{14}

\centerline{   \textbf{OR}}

Calculate Karl Pearson's coefficient of correlation from the following data:\pts{14}

\begin{center}
  
\begin{tabular}{lccccc}
     oprule
    X & 6 & 2 & 10 & 4 & 8 \\
    Y & 9 & 11 & 5 & 8 & 7 \\
    ottomrule
  \end{tabular}
\end{center}

\medskip\rule{\linewidth}{0.2pt}




\noindent   \textbf{Question 5.}
What is Regression Line? How it can be measured?\pts{14}

\centerline{   \textbf{OR}}

The following data based on 450 students are given for marks in Statistics and Economics at a certain examination. Find two regression lines and estimate the average marks in Economics of the candidates who obtained 50 marks in statistics:\pts{14}

\begin{center}
  
\begin{tabular}{lcc}
     oprule
    & Statistics & Economics \\
    \midrule
    Mean marks & 40 & 48 \\
    Standard deviation & 12 & 16 \\
    ottomrule
  \end{tabular}
\end{center}
Sum of products of deviations of marks from their respective means = 42,075.

\vfill
\rule{\linewidth}{0.4pt}

\begin{center}\small
\href{https://bhu-exam.vercel.app/}{Click here to visit our website}\\[4pt]
\href{https://me-aryan.vercel.app/}{Click here to visit our contributor}\\[4pt]
\href{https://quashan-me.vercel.app/}{Click here to visit our contributor}
\end{center}

\end{document}
